\subsection{ロピタルの定理の証明}
ロピタルの定理を証明しよう.微分係数の定義に帰着させる方法では,\,$c$ での微分係数だけを使った.今度は,\,$c$ と $x$ の間のどこかの点での導関数の比が必要になる.それを与えるのが,コーシーの平均値の定理(定理\ref{定理:コーシー平均値})であった.すなわち,区間 $[u,v]$ での増分の比は,ある $\xi\in(u,v)$ での導関数の比 $\frac{f'(\xi)}{g'(\xi)}$ に等しい.

証明は,\,$0$ と $0$ の形で端点が有限の場合,無限大と無限大の形の場合,\,$0$ と $0$ の形で端点が無限大の場合の順に進める.

\paragraphbf{$0$ と $0$ の形で端点が有限の場合}
有限の端点 $c$ で $0$ と $0$ の形になる場合を考える.\,$f(c)=g(c)=0$ と定めれば,\,$f,g$ は端点 $c$ まで含めて連続になるので,\,$c$ と $x$ を両端とする区間でコーシーの平均値の定理を使える.\,$f(x)-f(c)=f(x)$,\,$g(x)-g(c)=g(x)$ だから,
\[
 \frac{f(x)}{g(x)}=\frac{f'(\xi_x)}{g'(\xi_x)}
\]
となる点 $\xi_x$ が $c$ と $x$ の間にある.点 $\xi_x$ は $x$ によって変わるので,添字 $x$ を付けた.\,$x\to c$ なら,間に挟まれた $\xi_x$ も $c$ へ近づく.右辺は $L$ へ近づくから,左辺も $L$ へ近づく.

\paragraphbf{無限大と無限大の形の場合}
無限大と無限大の形では,もう一段の工夫が要る.\,$x\to+\infty$ で $|g(x)|\to\infty$,\,$L$ が実数の場合を示そう.ここでは,数列の極限の定義(\sectionref{節:数列極限})と同じように,誤差 $\varepsilon$ を使って極限を表す.すなわち,\,$t\to+\infty$ で $\frac{f'(t)}{g'(t)}\to L$ であるとは,どんな正の数 $\varepsilon$ に対しても,ある数 $B$ があって,\,$t>B$ なら
\[
 \left|\frac{f'(t)}{g'(t)}-L\right|<\varepsilon
\]
となることである.

正の数 $\varepsilon$ を任意にとり,上の $B$ を1つ選んで固定する.\,$x>B$ なら,\,$[B,x]$ でコーシーの平均値の定理を使うと
\[
 D_x=\frac{f(x)-f(B)}{g(x)-g(B)}=\frac{f'(\xi)}{g'(\xi)}
\]
となる $\xi\in(B,x)$ がある.\,$\xi>B$ だから $|D_x-L|<\varepsilon$ である.

$D_x$ は増分の比であり,求めたい商 $\frac{f(x)}{g(x)}$ そのものではない.そこで,元の商を $D_x$ で表してやる.\,$|g(x)|\to\infty$ なので,\,$x$ が十分大きければ $g(x)\ne0$ である.\,$D_x$ の定義から $f(x)=D_x\{g(x)-g(B)\}+f(B)$ となるので,両辺を $g(x)$ で割って
\[
 \frac{f(x)}{g(x)}
 =D_x\left(1-\frac{g(B)}{g(x)}\right)+\frac{f(B)}{g(x)}
\]
を得る.この両辺から $L=L\left(1-\frac{g(B)}{g(x)}\right)+\frac{Lg(B)}{g(x)}$ を引くと,
\[
 \frac{f(x)}{g(x)}-L
 =(D_x-L)\left(1-\frac{g(B)}{g(x)}\right)+\frac{f(B)-Lg(B)}{g(x)}
\]
となる.\,$B$ は固定してあるので,\,$f(B)$ と $g(B)$ は定数である.\,$|g(x)|\to\infty$ だから,ある $X>B$ があって,\,$x>X$ なら
\[
 \left|\frac{g(B)}{g(x)}\right|<1,\qquad
 \left|\frac{f(B)-Lg(B)}{g(x)}\right|<\varepsilon
\]
となる.1つ目の不等式から $\left|1-\frac{g(B)}{g(x)}\right|<2$ なので,\,$x>X$ のとき
\[
 \left|\frac{f(x)}{g(x)}-L\right|
 <\varepsilon\cdot2+\varepsilon=3\varepsilon
\]
である.\,$\varepsilon$ は任意だったので,\,$\frac{f(x)}{g(x)}\to L$ が示された.有限の端点で $|g|\to\infty$ となる場合も,\,「$t>B$」の代わりに「$t$ が端点から $\delta$ 以内」という範囲をとれば,同じ議論を使える.

\paragraphbf{$0$ と $0$ の形で端点が無限大の場合}
$0$ と $0$ の形で端点が $+\infty$ の場合は,同じ $B$ をとり,\,$B<x<y$ を満たす $x,y$ の間でコーシーの平均値の定理を使う.増分の比 $\frac{f(y)-f(x)}{g(y)-g(x)}$ は $(x,y)$ のどこかの点での $\frac{f'}{g'}$ に等しいので,\,$L$ との差は $\varepsilon$ 未満である.ここで $x$ を固定して $y\to+\infty$ とすると,\,$f(y)\to0$,\,$g(y)\to0$ より,この比は $\frac{-f(x)}{-g(x)}=\frac{f(x)}{g(x)}$ へ近づく.よって $x>B$ なら $\left|\frac{f(x)}{g(x)}-L\right|\le\varepsilon$ となり,\,$\frac{f(x)}{g(x)}\to L$ である.

残る場合も同じように示せる.\,$L=+\infty$ の場合は,「$L$ との差が $\varepsilon$ 未満」の代わりに「任意に与えた定数より大きい」を用い,\,$L=-\infty$ では不等号を逆にする.端点が $-\infty$ の場合は,\,「$t>B$」を「$t<B$」に替える.
